% !TeX program = pdflatex
\documentclass[11pt,a4paper]{article}
\usepackage[utf8]{inputenc}
\usepackage[T1]{fontenc}
\usepackage[brazil]{babel}
\usepackage{amsmath,amssymb,amsthm,amscd}
\usepackage[margin=2.4cm]{geometry}
\usepackage{microtype}
\usepackage[hidelinks]{hyperref}

\newtheorem{teorema}{Teorema}
\newtheorem{lema}[teorema]{Lema}
\newtheorem{proposicao}[teorema]{Proposição}
\newtheorem{corolario}[teorema]{Corolário}
\newtheorem{definicao}[teorema]{Definição}
\newtheorem{observacao}[teorema]{Observação}

\title{\textbf{A Cadeia Finita de Fermat, Canonicalização da Fase e Obstrução de Holonomia}}
\author{Aarão Melo Lopes}
\date{\today}

\begin{document}
\maketitle

\begin{abstract}
Para uma hipotética solução primitiva $a^k+b^k=c^k$, com $k>2$, a aplicação do Teorema de Zsigmondy a $c^k-b^k$ garante a existência de um divisor primo primitivo $p \mid (c^k-b^k)$ tal que a ordem multiplicativa de $u_p = bc^{-1}$ em $\mathbb{F}_p^\times$ é exatamente $k$. Fixada uma convenção de realização que associa o elemento $u_p$ de ordem $k$ à fase fundamental $2\pi/k$ na subálgebra real $\langle 1, J \rangle$ com $J^2 = -1$, define-se a representação $\widehat{u}_{p,k} = \exp(\frac{2\pi}{k}J)$, que induz o operador $A_0(k) = \operatorname{diag}(e^{2\pi i/k}, e^{-2\pi i/k}) \in \operatorname{SL}(2,\mathbb{C})$. Os operadores de curvatura do refinamento de Koch são definidos via \(\widetilde{J}=M_w^{-1}JM_w\) por
\[
Q_1=e^{\pi\widetilde{J}/6},\quad
Q_2=e^{-\pi\widetilde{J}/3},\quad
Q_3=e^{\pi\widetilde{J}/6},
\]
e \(A_1:=\Xi(A_0)=A_0Q_1A_0Q_2A_0Q_3A_0\).

\textbf{Teorema geométrico (fechado dentro do modelo).}
O desvio de traço fatora como \(D_{\mathrm{geom}}(c)=(1-c)P_D(c)\). Demonstra-se \(P_D(\cos(2\pi/k))\neq0\) para todo \(k>2\), donde \(\operatorname{tr}(A_1)\neq\operatorname{tr}(A_0)\). Para \(k=4\) obtém-se explicitamente \(\operatorname{tr}(A_1)=11\notin\{-2,0,2\}\), logo \(A_1^4\neq I\).

\textbf{Teorema aritmético necessário (ainda aberto).}
Pela Proposição de fatoração, \(\Xi\circ\mathcal{R}\) é constante nas fibras de \(\kappa\): \((\Xi\circ\mathcal{R})(a,b,c,p)=\widetilde{\Xi}(k)\). A composição atual é uma construção geométrica indexada por \(k\), não uma construção que ainda carregue \((a,b,c,p)\). Portanto
\[
\text{Zsigmondy + realização angular + Koch + obstrução espectral}\;\not\Rightarrow\;\text{FLT}.
\]
Fechar o Último Teorema de Fermat por esta via exigiria abandonar a fatoração \(\mathcal{R}=\widetilde{\mathcal{R}}\circ\kappa\) em favor de uma realização enriquecida que preserve informação aritmética suficiente até a entrada de \(\Xi\).
\end{abstract}

\section{Seleção Aritmética do Primo Primitivo}

\begin{lema}[Primo Primitivo de Zsigmondy]
Aplique o Teorema de Zsigmondy a $c^k-b^k$, com $c>b>0$ e $k>2$. Suponha uma solução primitiva $a^k+b^k=c^k$, com $\gcd(a,b,c)=1$. Então existe um primo $p \mid a$ tal que:
\begin{enumerate}
    \item $p \mid (c^k - b^k)$ e $p \nmid (c^m - b^m)$ para todo $1 \le m < k$;
    \item $p \nmid bc$;
    \item $\operatorname{ord}_{\mathbb{F}_p^\times}(bc^{-1}) = k$.
\end{enumerate}
Em particular, pelo Teorema de Lagrange, $k \mid (p-1)$, o que garante $p \nmid k$.
\end{lema}

\begin{proof}
Da primitividade $\gcd(a,b,c)=1$ segue $\gcd(a,b)=\gcd(a,c)=\gcd(b,c)=1$. De $a^k+b^k=c^k$ com $a,b>0$, temos automaticamente $c > b > 0$. Assim, a aplicação em $c^k - b^k = a^k$ está legitimada. O Teorema de Zsigmondy garante a existência de um divisor primo primitivo $p$ de $c^k - b^k$, exceto no caso $(c,b,k)=(2,1,6)$, o qual implicaria $a^6 = 2^6 - 1^6 = 63$, impossível pois $63$ não é uma sexta potência inteira.

Como $p \mid (c^k - b^k) = a^k$, segue $p \mid a$. Pela coprimalidade da solução ($\gcd(a,b)=\gcd(a,c)=1$), temos imediatamente $p \nmid b$ e $p \nmid c$, garantindo $p \nmid bc$. Como $p$ é divisor primitivo de $c^k-b^k$, $p \nmid (c^m - b^m)$ para $m < k$. Logo, $u_p = bc^{-1} \in \mathbb{F}_p^\times$ satisfaz $(bc^{-1})^k \equiv 1 \pmod p$. Se a ordem de $u_p$ fosse $d < k$, teríamos $p \mid (c^d - b^d)$, contradizendo a primitividade de $p$. Logo, o elemento possui ordem multiplicativa exata
\[
\operatorname{ord}_{\mathbb{F}_p^\times}(u_p) = k.
\]
Pelo Teorema de Lagrange, $k \mid (p-1) \implies p \nmid k$.
\end{proof}

\section{A Ponte Espectral e a Fatoração da Realização}

A estrutura cíclica abstrata de ordem $k$ estabelecida em $\mathbb{F}_p^\times$ admite uma representação complexa para análise de holonomia. A passagem não é canônica: trata-se de uma convenção de realização.

\begin{definicao}[Convenção de Realização Angular]
O elemento $u_p$ de ordem $k$ é associado, por convenção, à fase $2\pi/k$ na subálgebra real $\langle 1, J \rangle$ (com $J^2 = -1$). Representamos o ciclo gerado por $u_p$ pelo elemento:
\[
\widehat{u}_{p,k} := \exp\left(\frac{2\pi}{k}J\right) = \cos\frac{2\pi}{k}\cdot 1 + \sin\frac{2\pi}{k}\cdot J.
\]
Sob a leitura coordenada $1 \longleftrightarrow 1$ e $J \longleftrightarrow i$, a representação complexa do elemento de holonomia base assume a forma:
\[
A_0(k) = \begin{pmatrix} e^{2\pi i/k} & 0 \\ 0 & e^{-2\pi i/k} \end{pmatrix} \in \operatorname{SL}(2, \mathbb{C}),
\]
com $\operatorname{tr}(A_0) = 2\cos(2\pi/k)$ e $A_0(k)^k = I$.
\end{definicao}

\begin{proposicao}[Fatoração da Realização pela Ordem $k$]
Seja $\mathcal{S}$ o conjunto de tuplas $(a,b,c,p)$ correspondentes a soluções primitivas e ao seu primo de Zsigmondy associado. A aplicação de realização angular $\mathcal{R}:\mathcal{S}\to\operatorname{SL}(2,\mathbb{C})$ fatora através da projeção da ordem $\kappa(a,b,c,p)=\operatorname{ord}_{\mathbb{F}_p^\times}(bc^{-1})=k$, tal que
\[
\mathcal{R}=\widetilde{\mathcal{R}}\circ\kappa,\qquad\text{com }\widetilde{\mathcal{R}}(k)=A_0(k).
\]
Consequentemente,
\[
\Xi\circ\mathcal{R}=(\Xi\circ\widetilde{\mathcal{R}})\circ\kappa.
\]
Existe portanto uma única aplicação $\widetilde{\Xi}(k):=\Xi(A_0(k))$ tal que
\[
(\Xi\circ\mathcal{R})(a,b,c,p)=\widetilde{\Xi}(k).
\]
Em particular, $\Xi\circ\mathcal{R}$ é constante nas fibras de $\kappa$: se $\kappa(s_1)=\kappa(s_2)$, então $(\Xi\circ\mathcal{R})(s_1)=(\Xi\circ\mathcal{R})(s_2)$. O objeto produzido pelo pipeline atual não distingue duas quádruplas com o mesmo $k$.
\end{proposicao}

\begin{observacao}[Status da convenção]
A associação $u_p\mapsto\widehat{u}_{p,k}$ não é um homomorfismo ou functor canônico. Trata-se de uma escolha de realização angular. A Proposição acima transforma a perda de informação de $(a,b,c,p)$ em uma propriedade matemática da aplicação: $\Xi\circ\mathcal{R}$ depende exclusivamente de $k$.
\end{observacao}

\section{Interface de Viviani e Classificação Espectral}

No Corpo Universal, a curva de Viviani (interseção da esfera \(x^2+y^2+z^2=4a^2\) com o cilindro tangente \((x-a)^2+y^2=a^2\)) realiza o regime associado ao parâmetro \(d=+1\) da dobra de renormalização. A recorrência
\[
t_k = m\, t_{k-1} - d\, t_{k-2}
\]
distingue dois casos internos:
\begin{center}
\begin{tabular}{c|c|c}
\(d\) & interpretação no corpus & \\
\hline
\(-1\) & hiperbólico / rectilíneo & Cantor--Julia (corte) \\
\(+1\) & elíptico / rotação & \textbf{Viviani} \\
\end{tabular}
\end{center}
Quando \(d=+1\) a dinâmica roda (meio-ângulo); quando \(d=-1\) ela cresce (corte). A interface de Viviani é o lugar geométrico natural onde se situa a holonomia inicial \(A_0\) (rotação de ângulo \(2\pi/k\)).

Para conectar essa estrutura interna à classificação clássica de elementos de \(\operatorname{SL}(2,\mathbb{C})\), recordamos que se \(A\in\operatorname{SL}(2,\mathbb{C})\) então o polinômio característico é \(\chi_A(\lambda)=\lambda^2-\operatorname{tr}(A)\lambda+1\) e o discriminante espectral é
\[
\Delta_A=\operatorname{tr}(A)^2-4.
\]
No caso real:
\begin{align*}
\lvert\operatorname{tr}(A)\rvert<2 &\quad\Longrightarrow\quad\text{elíptico},\\
\lvert\operatorname{tr}(A)\rvert=2 &\quad\Longrightarrow\quad\text{parabólico/degenerado},\\
\lvert\operatorname{tr}(A)\rvert>2 &\quad\Longrightarrow\quad\text{hiperbólico}.
\end{align*}
A identificação entre o parâmetro \(d\) do corpus e o sinal de \(\Delta_A\) (ou, equivalentemente, a comparação de \(\lvert\operatorname{tr}(A)\rvert\) com \(2\)) é uma proposição estrutural que deve ser enunciada explicitamente quando se invoca a interface de Viviani em argumentos que cruzam os dois regimes. No presente texto, \(A_0\) satisfaz \(\lvert\operatorname{tr}(A_0)\rvert\le 2\) por construção (rotação), logo situa-se no regime elíptico clássico.

Propriedades operacionais medidas no anel (Teor.~\emph{viviani} do corpus):
\begin{itemize}
    \item a dobra é a lei: \(2\cos(2u)=(2\cos u)^2-2\) realiza o operador de renormalização \(R\) com membrana \(d=1\);
    \item o recobrimento é o espelho: troca de folha \((y,z)\mapsto(y,-z)\) com \(\mathrm{espelho}^2=I\);
    \item o nó é o ponto fixo \(R(2)=2\) da cascata;
    \item as projeções são os cortes (parábola e quártica de Gerono).
\end{itemize}

\section{Geometria do Operador de Refinamento}

Seja \(J\) o gerador da subálgebra real com \(J^2=-I\). Fixada uma conjugação \(M_w\), defina
\[
\widetilde{J}:=M_w^{-1}JM_w.
\]
Então \(\widetilde{J}^2=-I\). Os operadores de curvatura do refinamento de Koch são as exponenciais do mesmo gerador:
\[
Q_1=e^{\frac{\pi}{6}\widetilde{J}},\qquad
Q_2=e^{-\frac{\pi}{3}\widetilde{J}},\qquad
Q_3=e^{\frac{\pi}{6}\widetilde{J}}.
\]
Como são funções do mesmo operador, comutam, e
\[
Q_1Q_2Q_3=e^{(\pi/6-\pi/3+\pi/6)\widetilde{J}}=I.
\]
O refinamento é
\[
\Xi(H)=H\,Q_1\,H\,Q_2\,H\,Q_3\,H,
\]
e \(A_1:=\Xi(A_0)\).

Na família de holonomias considerada, a realização diagonal \(A_0\) satisfaz a parametrização de invariantes
\[
x=\operatorname{tr}(A_0)=2c,\qquad y=\operatorname{tr}(A_0\widetilde{J}),
\]
com \(c=\cos(2\pi/k)\). A aplicação \(\Xi(A_0)=A_1\) induz o desvio de traço bivariado \(D(x,y)=\operatorname{tr}(A_1)-\operatorname{tr}(A_0)\). A substituição das coordenadas geométricas resulta na fatoração univariada
\[
D_{\mathrm{geom}}(c):=D\bigl(2c,-4(1-c^2)\bigr)=(1-c)P_D(c),
\]
onde \(P_D(c)\in\mathbb{Z}[\sqrt{3}][c]\) é o polinômio explícito de grau 6:
\[
\begin{aligned}
P_D(c)={}&16\sqrt{3}\,c^6+(40+16\sqrt{3})c^5+(40-28\sqrt{3})c^4\\
&-(58+28\sqrt{3})c^3+(-58+12\sqrt{3})c^2+(9+12\sqrt{3})c+11.
\end{aligned}
\]
Como \(1-c_k>0\) para todo \(k>2\), a condição de ruptura \(\operatorname{tr}(A_1)\neq\operatorname{tr}(A_0)\) equivale a \(P_D(c_k)\neq0\).

O desvio \(D_{\mathrm{geom}}\) mede a variação do traço produzida pelo refinamento relativamente à holonomia inicial:
\[
D_{\mathrm{geom}}(c)=\operatorname{tr}(A_1)-\operatorname{tr}(A_0).
\]
Em particular, quando \(\lvert\operatorname{tr}(A_1)\rvert>2\), essa variação corresponde à saída do regime elíptico clássico para o regime hiperbólico. Essa saída é verificada para \(k=4\) (\(\operatorname{tr}(A_1)=11\)), mas não foi demonstrada para todo \(k>2\).

\section{Incompatibilidade Rigorosa para \texorpdfstring{$k=4$}{k=4}}

\begin{teorema}[Quebra Efetiva de Ordem Finita para $k=4$]
Para $k=4$, a convenção angular fixa
\[
A_0=\begin{pmatrix}i&0\\0&-i\end{pmatrix},\qquad\operatorname{tr}(A_0)=0.
\]
Com a conjugação que produz
\[
\widetilde{J}=\begin{pmatrix}2i&-1\\-3&-2i\end{pmatrix},
\]
obtém-se
\[
A_0\widetilde{J}=\begin{pmatrix}-2&-i\\3i&-2\end{pmatrix},\qquad
y=\operatorname{tr}(A_0\widetilde{J})=-4.
\]
A fórmula do desvio bivariado fornece então
\[
D(0,-4)=\frac18\bigl(6\cdot16-8\bigr)=\frac{88}{8}=11,
\]
logo
\[
\operatorname{tr}(A_1)=11.
\]
O polinômio característico de \(A_1\) é \(\lambda^2-11\lambda+1\), com raízes
\[
\lambda_\pm=\frac{11\pm\sqrt{117}}2
\]
(\(\lambda_+>1\), \(0<\lambda_-<1\)). Em particular \(\lvert\operatorname{tr}(A_1)\rvert=11>2\), de modo que \(A_1\) é hiperbólico. Como \(A^4=I\) implica necessariamente \(\operatorname{tr}(A)\in\{-2,0,2\}\) e \(11\notin\{-2,0,2\}\), conclui-se
\[
A_1^4\neq I.
\]
Assim a cadeia geométrica produz, para \(k=4\),
\[
A_0^4=I\qquad\text{mas}\qquad A_1^4\neq I.
\]
\end{teorema}

\section{Certificado Algorítmico de Irredutibilidade sobre \texorpdfstring{$K$}{K}}

\begin{lema}[Irredutibilidade do Polinômio Geométrico $P_D$]
O polinômio geométrico $P_D(c) \in \mathbb{Z}[\sqrt{3}][c]$ é irredutível sobre $K = \mathbb{Q}(\sqrt{3})$.
\end{lema}

\begin{proof}
Considere o ideal primo $\mathfrak{p} = (23, \sqrt{3} - 7) \subset \mathbb{Z}[\sqrt{3}]$, amparado por $7^2 = 49 \equiv 3 \pmod{23}$.
Sob a redução $\sqrt{3} \mapsto 7 \pmod{23}$, obtemos:
\[
\bar{P}_{23}(c) = -3c^6 - 9c^5 + 5c^4 - c^3 + 3c^2 + c + 11 \in \mathbb{F}_{23}[c].
\]
Como $-3$ é invertível em $\mathbb{F}_{23}$, normalizamos dividindo pelo coeficiente líder para obter o polinômio mônico:
\[
f(c) = c^6 + 3c^5 + 6c^4 + 8c^3 - c^2 - 8c + 4 \in \mathbb{F}_{23}[c].
\]
Os testes de Frobenius, no sentido do critério padrão de irredutibilidade modular, exigem
\[
x^{23^6}\equiv x\pmod{f}
\]
e, para cada primo $\ell\mid 6$ (isto é, $\ell=2$ e $\ell=3$),
\[
\gcd\bigl(f,x^{23^{6/\ell}}-x\bigr)=1.
\]
Explicitamente:
\[
\gcd\bigl(f,x^{23^3}-x\bigr)=1, \qquad \gcd\bigl(f,x^{23^2}-x\bigr)=1.
\]
Essas três condições (certificadas computacionalmente) excluem fatores irredutíveis de graus $2$ e $3$, enquanto a congruência de grau $6$ garante que todos os fatores têm grau divisor de $6$. Logo $f$ é irredutível de grau $6$ em $\mathbb{F}_{23}[c]$. Pelo critério de redução modular, $P_D(c)$ é irredutível sobre $K = \mathbb{Q}(\sqrt{3})$.
\end{proof}

\section{Não-Anulação de \texorpdfstring{$P_D$}{PD} e Ruptura de Classe de Conjugação}

\begin{teorema}[Ruptura Geral da Classe de Conjugação]
Para todo inteiro $k>2$, vale
\[
\operatorname{tr}(A_1)\neq\operatorname{tr}(A_0).
\]
Equivalentemente,
\[
P_D(c_k)\neq0,\qquad c_k=\cos\frac{2\pi}{k}.
\]
\end{teorema}

\begin{proof}
Suponha, por absurdo, que exista $k>2$ tal que
\[
P_D(c_k)=0.
\]
Pela irredutibilidade de $P_D$ sobre
\[
K=\mathbb{Q}(\sqrt{3}),
\]
e pelo fato de $\deg P_D=6$, segue que $P_D$ é o polinômio mínimo de $c_k$ sobre $K$. Portanto,
\[
[K(c_k):K]=6.
\]
Defina
\[
F_k=\mathbb{Q}(c_k).
\]
Como $K/\mathbb{Q}$ é uma extensão quadrática de Galois, vale a fórmula de graus
\[
[KF_k:K]=[F_k:K\cap F_k].
\]
Assim,
\[
6=[F_k:K\cap F_k].
\]
Como
\[
[K\cap F_k:\mathbb{Q}]\in\{1,2\},
\]
obtemos
\[
[F_k:\mathbb{Q}]=[F_k:K\cap F_k]\,[K\cap F_k:\mathbb{Q}]\in\{6,12\}.
\]
Para $k>2$,
\[
[F_k:\mathbb{Q}]=\left[\mathbb{Q}\!\left(\cos\frac{2\pi}{k}\right):\mathbb{Q}\right]=\frac{\varphi(k)}{2}.
\]
Consequentemente,
\[
\varphi(k)\in\{12,24\}.
\]
Os únicos inteiros $k>2$ com $\varphi(k)\in\{12,24\}$ são
\[
\mathcal{K}=\{13,21,26,28,35,36,39,42,45,52,56,70,72,78,84,90\}.
\]
Resta, portanto, excluir apenas esses casos.

Considere o ideal primo
\[
\mathfrak{q}=(11,\sqrt{3}-5)\subset\mathbb{Z}[\sqrt{3}],
\]
que é primo porque
\[
5^2\equiv3\pmod{11}.
\]
A redução de $P_D$ módulo $\mathfrak{q}$ é
\[
\overline{P}_D(c)=3c^6-c^5-c^4+2c^2+3c\in\mathbb{F}_{11}[c].
\]
Para cada $k\in\mathcal{K}$, seja $m_k(c)\in\mathbb{Z}[c]$ o polinômio mínimo de
$c_k=\cos(2\pi/k)$ sobre $\mathbb{Q}$, e defina
\[
R_k=\operatorname{Res}_c(P_D,m_k)\in\mathbb{Z}[\sqrt{3}].
\]
A redução dos resultantes no corpo residual
\[
\mathbb{Z}[\sqrt{3}]/\mathfrak{q}\simeq\mathbb{F}_{11}
\]
fornece:
\[
\begin{array}{c|c@{\qquad}c|c}
k&R_k\bmod\mathfrak{q}&k&R_k\bmod\mathfrak{q}\\
\hline
13&2&45&3\\
21&9&52&8\\
26&9&56&1\\
28&6&70&6\\
35&5&72&6\\
36&1&78&2\\
39&3&84&6\\
42&9&90&1
\end{array}
\]
e, portanto,
\[
R_k\bmod\mathfrak{q}\neq0
\]
para todo $k\in\mathcal{K}$.
Logo
\[
R_k\neq0
\]
em $K$. Assim,
\[
\gcd_{K[c]}(P_D,m_k)=1,
\]
o que implica
\[
P_D(c_k)\neq0
\]
para todo $k\in\mathcal{K}$.

Isto contradiz a hipótese inicial $P_D(c_k)=0$. Portanto,
\[
P_D(c_k)\neq0\qquad\text{para todo }k>2.
\]
Como
\[
1-c_k>0\qquad(k>2),
\]
da fatoração
\[
D_{\mathrm{geom}}(c_k)=(1-c_k)P_D(c_k)
\]
segue
\[
D_{\mathrm{geom}}(c_k)\neq0.
\]
Finalmente,
\[
\operatorname{tr}(A_1)-\operatorname{tr}(A_0)=D_{\mathrm{geom}}(c_k)\neq0,
\]
e portanto
\[
\operatorname{tr}(A_1)\neq\operatorname{tr}(A_0)
\]
para todo $k>2$.
\end{proof}

\section{Síntese da Fronteira Lógica}

\begin{observacao}[Dois níveis, status distinto]
A cadeia documentada no manuscrito é
\[
a^k+b^k=c^k
\ \Longrightarrow\
\operatorname{ord}_{\mathbb{F}_p^\times}(bc^{-1})=k
\ \Longrightarrow\
A_0^k=I
\ \Longrightarrow\
A_1=\Xi(A_0)
\]
e, independentemente,
\[
A_1=\Xi(A_0)
\ \Longrightarrow\
\operatorname{tr}(A_1)\neq\operatorname{tr}(A_0).
\]
Isso está efetivamente fechado. Em particular, para \(k=4\),
\[
A_0^4=I,\qquad\operatorname{tr}(A_1)=11,\qquad A_1^4\neq I.
\]
O caso \(k=4\) não é um ``contraexemplo teórico definitivo'' à hipótese de Fermat; ele contradiz qualquer princípio adicional do modelo que imponha \(A_1^4=I\).

A identidade \(Q_1Q_2Q_3=I\) \emph{não} permite reagrupar \(\Xi(A_0)\) como potência de \(A_0\), porque em geral \(Q_iA_0\neq A_0Q_i\). O próprio cálculo \(k=4\) demonstra a não-preservação.

Além disso, \(\operatorname{tr}(A_1)\neq\operatorname{tr}(A_0)\) não implica \(A_1^k\neq I\). Para contradizer \(A_1^k=I\) seria necessário
\[
\operatorname{tr}(A_1)\notin\mathcal{T}_k=\Bigl\{2\cos\frac{2\pi j}{k}:0\le j<k\Bigr\}.
\]
Isso vale para \(k=4\) (\(\mathcal{T}_4=\{-2,0,2\}\), \(\operatorname{tr}(A_1)=11\)), mas não foi demonstrado para \(k>4\).

\begin{center}
\begin{tabular}{@{}lp{9.0cm}@{}}
\hline
Componente & Status \\
\hline
Zsigmondy $\to\operatorname{ord}(u_p)=k$ & fechado \\
Proposição de fatoração $\mathcal{R}=\widetilde{\mathcal{R}}\circ\kappa$ & fechado \\
Invariância $Q_1Q_2Q_3=I$ & fechado \\
Não-anulação $P_D(c_k)\neq0$ & fechado \\
$k=4$: $\operatorname{tr}(A_1)=11\Rightarrow A_1^4\neq I$ & fechado \\
Preservação genérica $A_1^k=I$ & incompatível em $k=4$ \\
Ponte aritmética enriquecida (não fatorada por $k$) & aberta \\
\hline
\end{tabular}
\end{center}

O ponto estrutural mais forte é o seguinte. A composição atual \(\Xi\circ\mathcal{R}\) é uma construção geométrica indexada por \(k\), e não uma construção que ainda carregue \((a,b,c,p)\). Formalmente o esquema é
\[
(a,b,c,p)\longrightarrow u_p\longrightarrow k\longrightarrow A_0(k)\longrightarrow A_1(k).
\]
Logo \(A_1=A_1(k)\): não existe dependência explícita de \(a,b,c,p\) que possa posteriormente ser usada para deduzir uma restrição adicional sobre \(A_1\). Após a convenção angular \(u_p\rightsquigarrow e^{2\pi i/k}\), dois dados aritméticos completamente diferentes que produzam o mesmo \(k\) chegam ao mesmo \(A_0\) e, portanto, ao mesmo \(A_1\).

Consequência metodológica: para fechar o Último Teorema de Fermat por esta via seria necessário modificar a própria realização, de modo que o dado aritmético sobreviva até a entrada de \(\Xi\). Idealmente, um diagrama genuinamente comutativo
\[
\begin{CD}
\text{dados de Fermat}@>{\mathcal{A}}>>\text{objeto geométrico enriquecido}\\
@V{\text{aritmética}}VV@VV{\Xi}V\\
\text{invariante aritmético}@>{\Phi}>>\text{invariante após refinamento}
\end{CD}
\]
com uma implicação
\[
a^k+b^k=c^k\quad\Longrightarrow\quad\Phi(\Xi(A))=\Phi(A),
\]
enquanto o cálculo geométrico demonstra \(\Phi(\Xi(A))\neq\Phi(A)\). Na ausência desse diagrama (ou de estrutura matematicamente equivalente),
\[
\text{Zsigmondy + realização angular + Koch + obstrução espectral}\;\not\Rightarrow\;\text{FLT}.
\]
O caso \(k=4\) mostra ainda que não se pode simplesmente postular a ponte depois do fato: qualquer lema genérico afirmando que ``a origem Fermat implica preservação da ordem sob \(\Xi\)'' entraria em conflito com a dinâmica geométrica já calculada.

Neste estágio o polinômio \(P_D\) e a obstrução espectral dentro do modelo consideram-se encerrados. A próxima etapa é uma investigação independente:
\begin{quote}
Existe uma realização aritmético-geométrica não degenerada que preserve informação suficiente da equação de Fermat para produzir uma restrição sobre a holonomia refinada?
\end{quote}
Se a resposta for negativa, a própria fatoração da arquitetura atual por \(k\) torna-se um resultado matemático: a construção não distingue soluções aritméticas com o mesmo expoente.
\end{observacao}

\end{document}
